% ============================================================
%  CHAPITRE III : Conception et réalisation
% ============================================================
\chapter{Conception et réalisation}
\label{chap:conception}

\section*{Introduction}

Ce chapitre détaille la conception architecturale du système proposé, les choix technologiques ainsi que les étapes de mise en œuvre. Nous présentons également les diagrammes \gls{uml} et l'architecture globale de la solution.

% -------------------------------------------------------
\section{Architecture globale du système}
\label{sec:architecture}

\subsection{Vue d'ensemble}

La figure~\ref{fig:architecture} illustre l'architecture globale du système de détection d'intrusions proposé.

\begin{figure}[H]
  \centering
  \begin{tikzpicture}[
    node distance=1.2cm and 2cm,
    block/.style={
      rectangle, draw=iacsTeal, fill=iacsLight!50,
      rounded corners=5pt, minimum width=3cm,
      minimum height=1.2cm, align=center,
      font=\small\sffamily, line width=0.8pt,
      drop shadow={shadow xshift=1pt, shadow yshift=-1pt, opacity=0.1}
    },
    dbblock/.style={
      cylinder, draw=iacsTeal, fill=iacsGreen!15,
      shape border rotate=90, aspect=0.25,
      minimum width=2cm, minimum height=1.2cm,
      font=\small\sffamily, line width=0.8pt
    },
    alertblock/.style={
      block, fill=iacsOrange!15, draw=iacsOrange
    },
    arrow/.style={-Stealth, thick, color=iacsTeal},
    dataarrow/.style={-Stealth, thick, color=iacsAccent, dashed}
  ]
    % Couche collecte
    \node[block] (capture) {Capture\\trafic réseau};
    \node[block, right=of capture] (preprocess) {Pré-traitement\\des données};
    \node[block, right=of preprocess] (extract) {Extraction de\\caractéristiques};

    % Couche IA
    \node[block, below=1.5cm of extract, fill=iacsTeal!15] (model) {\textbf{Modèle CNN}\\classification};
    \node[dbblock, left=of model] (db) {Base de\\données};

    % Couche décision
    \node[alertblock, below=1.5cm of model] (alert) {\faExclamationTriangle\; Système\\d'alerte};
    \node[block, left=of alert] (dashboard) {\faDesktop\; Dashboard\\monitoring};

    % Flèches
    \draw[arrow] (capture) -- (preprocess);
    \draw[arrow] (preprocess) -- (extract);
    \draw[arrow] (extract) -- (model);
    \draw[dataarrow] (model) -- (db);
    \draw[arrow] (model) -- (alert);
    \draw[dataarrow] (alert) -- (dashboard);
    \draw[dataarrow] (db) -- (dashboard);

    % Labels couches
    \node[font=\footnotesize\sffamily\bfseries\color{iacsTeal},
          anchor=west] at ([xshift=-1cm, yshift=0.8cm]capture.north) 
          {\faLayerGroup\; Couche collecte};
    \node[font=\footnotesize\sffamily\bfseries\color{iacsTeal},
          anchor=east] at ([xshift=1cm, yshift=0.8cm]model.north) 
          {\faBrain\; Couche IA};
    \node[font=\footnotesize\sffamily\bfseries\color{iacsOrange},
          anchor=east] at ([xshift=1cm, yshift=0.8cm]alert.north) 
          {\faShieldAlt\; Couche décision};
  \end{tikzpicture}
  \caption{Architecture globale du système de détection d'intrusions}
  \label{fig:architecture}
\end{figure}

% -------------------------------------------------------
\section{Conception détaillée}
\label{sec:conception_detail}

\subsection{Diagramme de cas d'utilisation}

\hspace{1cm}Le diagramme de cas d'utilisation (figure~\ref{fig:use_case}) présente les interactions principales entre les acteurs et le système.

\begin{figure}[H]
  \centering
  \begin{tikzpicture}[
    actor/.style={font=\small\sffamily},
    usecase/.style={
      ellipse, draw=iacsTeal, fill=iacsLight!30,
      minimum width=3cm, minimum height=1cm,
      font=\small\sffamily, align=center
    }
  ]
    % Acteurs
    \node[actor] (admin) at (0,0) {\faUserShield};
    \node[below=0.1cm of admin, font=\scriptsize\sffamily] {Administrateur};

    \node[actor] (analyst) at (0,-4) {\faUserCog};
    \node[below=0.1cm of analyst, font=\scriptsize\sffamily] {Analyste SOC};

    % Cas d'utilisation
    \node[usecase] (monitor) at (5, 0) {Monitorer\\le trafic};
    \node[usecase] (config) at (5, -1.5) {Configurer\\les règles};
    \node[usecase] (analyze) at (5, -3) {Analyser\\les alertes};
    \node[usecase] (report) at (5, -4.5) {Générer\\des rapports};

    % Relations
    \draw[iacsTeal] (admin) -- (monitor);
    \draw[iacsTeal] (admin) -- (config);
    \draw[iacsTeal] (analyst) -- (analyze);
    \draw[iacsTeal] (analyst) -- (report);
    \draw[iacsTeal, dashed] (admin) -- (analyze);

    % Système
    \draw[iacsTeal, rounded corners=8pt, thick] 
      (3, 1) rectangle (7.5, -5.5);
    \node[font=\footnotesize\sffamily\bfseries\color{iacsTeal}] 
      at (5.25, 0.7) {Système IDS-IA};
  \end{tikzpicture}
  \caption{Diagramme de cas d'utilisation}
  \label{fig:use_case}
\end{figure}

\subsection{Architecture du modèle CNN}

Le réseau proposé se compose des couches suivantes (figure~\ref{fig:cnn_arch}) :

\begin{figure}[H]
  \centering
  \begin{tikzpicture}[
    layer/.style={
      rectangle, rounded corners=3pt,
      minimum width=1.8cm, minimum height=0.8cm,
      font=\scriptsize\sffamily\bfseries, align=center,
      line width=0.5pt
    },
    arrow/.style={-Stealth, thick, color=iacsGray}
  ]
    \node[layer, draw=iacsTeal, fill=iacsTeal!20] (input) {Input\\$n \times m$};
    \node[layer, draw=iacsAccent, fill=iacsAccent!15, right=0.5cm of input] (conv1) {Conv1D\\64 filtres};
    \node[layer, draw=iacsAccent, fill=iacsAccent!15, right=0.5cm of conv1] (pool1) {MaxPool\\$2 \times 1$};
    \node[layer, draw=iacsGreen!70!black, fill=iacsGreen!15, right=0.5cm of pool1] (conv2) {Conv1D\\128 filtres};
    \node[layer, draw=iacsGreen!70!black, fill=iacsGreen!15, right=0.5cm of conv2] (pool2) {MaxPool\\$2 \times 1$};
    \node[layer, draw=iacsOrange, fill=iacsOrange!15, right=0.5cm of pool2] (fc) {Dense\\256};
    \node[layer, draw=iacsRed!70, fill=iacsRed!10, right=0.5cm of fc] (out) {Softmax\\$k$ classes};

    \draw[arrow] (input) -- (conv1);
    \draw[arrow] (conv1) -- (pool1);
    \draw[arrow] (pool1) -- (conv2);
    \draw[arrow] (conv2) -- (pool2);
    \draw[arrow] (pool2) -- (fc);
    \draw[arrow] (fc) -- (out);
  \end{tikzpicture}
  \caption{Architecture du modèle CNN proposé}
  \label{fig:cnn_arch}
\end{figure}

% -------------------------------------------------------
\section{Environnement et outils de développement}
\label{sec:outils}

\subsection{Technologies utilisées}

Le tableau~\ref{tab:technologies} résume les technologies et frameworks utilisés.

\begin{table}[H]
  \centering
  \caption{Technologies et outils de développement}
  \label{tab:technologies}
  \renewcommand{\arraystretch}{1.4}
  \begin{tabular}{|>{\centering\arraybackslash}m{2.5cm}|l|l|}
    \hline
    \rowcolor{iacsTableHead}
    \textcolor{white}{\textbf{Catégorie}} & 
    \textcolor{white}{\textbf{Outil}} & 
    \textcolor{white}{\textbf{Version}} \\
    \hline
    \multirow{2}{*}{Langage} 
      & Python & 3.11 \\
      \cline{2-3}
      & JavaScript & ES2023 \\
    \hline
    \rowcolor{iacsLightGray}
    \multirow{2}{*}{Frameworks IA} 
      & TensorFlow & 2.15 \\
      \cline{2-3}
      \rowcolor{iacsLightGray}
      & Scikit-learn & 1.4 \\
    \hline
    \multirow{2}{*}{Sécurité} 
      & Scapy & 2.5 \\
      \cline{2-3}
      & Snort & 3.1 \\
    \hline
    \rowcolor{iacsLightGray}
    Base de données & MongoDB & 7.0 \\
    \hline
    Conteneurisation & Docker & 24.0 \\
    \hline
  \end{tabular}
\end{table}

\subsection{Captures d'écran}

Les figures~\ref{fig:screenshots} montrent des captures d'écran de l'interface développée.

\begin{figure}[H]
  \centering
  \begin{subfigure}[b]{0.45\textwidth}
    \centering
    \fcolorbox{iacsTeal}{iacsLight!20}{%
      \begin{minipage}{0.9\textwidth}
        \centering
        \vspace{1.5cm}
        {\small\sffamily\color{iacsGray} Capture d'écran 1\\
        \textit{(insérer image)}}
        \vspace{1.5cm}
      \end{minipage}
    }
    \caption{Interface de monitoring}
    \label{fig:screen_monitor}
  \end{subfigure}
  \hfill
  \begin{subfigure}[b]{0.45\textwidth}
    \centering
    \fcolorbox{iacsTeal}{iacsLight!20}{%
      \begin{minipage}{0.9\textwidth}
        \centering
        \vspace{1.5cm}
        {\small\sffamily\color{iacsGray} Capture d'écran 2\\
        \textit{(insérer image)}}
        \vspace{1.5cm}
      \end{minipage}
    }
    \caption{Tableau de bord des alertes}
    \label{fig:screen_alerts}
  \end{subfigure}
  \caption{Captures d'écran de l'interface}
  \label{fig:screenshots}
\end{figure}

\iacswarning{N'oubliez pas d'insérer vos propres captures d'écran en remplaçant les placeholders ci-dessus par~: \texttt{\textbackslash includegraphics[width=\textbackslash linewidth]\{figures/votre\_image.png\}}}

% -------------------------------------------------------
\section{Résultats et discussion}
\label{sec:resultats}

\subsection{Matrice de confusion}

La figure~\ref{fig:confusion} présente la matrice de confusion obtenue.

\begin{figure}[H]
  \centering
  \begin{tikzpicture}[scale=1.2]
    % Matrice
    \fill[iacsGreen!25] (0,1) rectangle (1,2);
    \fill[iacsRed!15] (1,1) rectangle (2,2);
    \fill[iacsRed!15] (0,0) rectangle (1,1);
    \fill[iacsGreen!25] (1,0) rectangle (2,2);
    \fill[iacsGreen!25] (1,0) rectangle (2,1);

    \draw[iacsTeal, thick] (0,0) grid (2,2);

    % Valeurs
    \node[font=\large\bfseries\color{iacsGreen!70!black}] at (0.5,1.5) {1842};
    \node[font=\large\bfseries\color{iacsRed}] at (1.5,1.5) {23};
    \node[font=\large\bfseries\color{iacsRed}] at (0.5,0.5) {31};
    \node[font=\large\bfseries\color{iacsGreen!70!black}] at (1.5,0.5) {1604};

    % Labels
    \node[font=\small\sffamily, anchor=south] at (0.5,2.1) {Normal};
    \node[font=\small\sffamily, anchor=south] at (1.5,2.1) {Attaque};
    \node[font=\small\sffamily\bfseries, anchor=south] at (1,2.5) {Prédit};

    \node[font=\small\sffamily, rotate=90, anchor=south] at (-0.3,1.5) {Normal};
    \node[font=\small\sffamily, rotate=90, anchor=south] at (-0.3,0.5) {Attaque};
    \node[font=\small\sffamily\bfseries, rotate=90, anchor=south] at (-0.7,1) {Réel};
  \end{tikzpicture}
  \caption{Matrice de confusion du modèle proposé}
  \label{fig:confusion}
\end{figure}

\subsection{Résultats quantitatifs}

Le tableau~\ref{tab:resultats} résume les performances obtenues.

\begin{table}[H]
  \centering
  \caption{Résultats de classification par type d'attaque}
  \label{tab:resultats}
  \renewcommand{\arraystretch}{1.3}
  \begin{tabular}{|l|c|c|c|c|}
    \hline
    \rowcolor{iacsTableHead}
    \textcolor{white}{\textbf{Classe}} & 
    \textcolor{white}{\textbf{Précision}} & 
    \textcolor{white}{\textbf{Rappel}} & 
    \textcolor{white}{\textbf{F1-score}} & 
    \textcolor{white}{\textbf{Support}} \\
    \hline
    Normal    & 0,983 & 0,988 & 0,985 & 1\,865 \\
    \hline
    \rowcolor{iacsLightGray}
    DoS       & 0,972 & 0,965 & 0,968 & 892 \\
    \hline
    Probe     & 0,951 & 0,943 & 0,947 & 421 \\
    \hline
    \rowcolor{iacsLightGray}
    R2L       & 0,934 & 0,912 & 0,923 & 215 \\
    \hline
    U2R       & 0,918 & 0,897 & 0,907 & 107 \\
    \hline
    \rowcolor{iacsTeal!10}
    \textbf{Moyenne} & \textbf{0,963} & \textbf{0,958} & \textbf{0,960} & \textbf{3\,500} \\
    \hline
  \end{tabular}
\end{table}

\section*{Conclusion}

Ce chapitre a présenté la conception et la réalisation du système proposé. Les résultats obtenus démontrent l'efficacité de l'approche basée sur les \gls{cnn} pour la détection d'intrusions, avec une précision globale de 96,3\%.
